%! Absolute Value Functions \& Transformations — Unit 2, Lesson 2.3 — Lesson Plan
\documentclass[10pt]{article}
\usepackage{algebra2-article}
\usepackage{algebra2-boxes}
\usepackage{graphicx}   % -article does NOT load graphicx; needed for thumbnails/screenshots
\graphicspath{{images/}}
\newcommand{\TallMath}[1]{$\displaystyle #1\rule[-1.4em]{0pt}{3.2em}$}
\newcommand{\CourseName}{Algebra 2: Shepherd}
\newcommand{\SchoolYear}{2026--2027}
\newcommand{\MeetingLength}{55 minutes}
\newcommand{\UnitNumberName}{Unit 2: Linear Functions \quad}
\newcommand{\LessonNumberName}{Lesson 2.3: Absolute Value Functions \& Transformations}
% Coordinate grid: {xmin}{xmax}{ymin}{ymax}
\newcommand{\lgrid}[4]{%
  \draw[step=1, linegray!40, very thin] (#1,#3) grid (#2,#4);
  \draw[->, semithick, charcoal] (#1,0)--(#2,0) node[below right, font=\scriptsize]{$x$};
  \draw[->, semithick, charcoal] (0,#3)--(0,#4) node[left, font=\scriptsize]{$y$};}

\begin{document}
\begin{center}
    {\Large\bfseries \CourseName: \SchoolYear \\
    {\small\bfseries \UnitNumberName \LessonNumberName}}
\end{center}

\begin{tcolorbox}[colback=forestbg, colframe=forest, boxrule=0.9pt, arc=2mm,
    left=2mm, right=2mm, top=1.5mm, bottom=1.5mm]
    \textbf{Primary Objective:} Students can read the absolute value parent function $f(x)=|x|$ as the
    graph of \emph{distance from zero} --- a \textbf{V-shape} --- and describe every other absolute
    value function as a \textbf{transformation} of that parent. From the vertex form
    $g(x)=a|x-h|+k$ they identify the \emph{vertex} $(h,k)$, the \emph{axis of symmetry} $x=h$, the
    direction of opening (up if $a>0$, down if $a<0$), the \emph{minimum or maximum} value $k$, and
    whether the graph is \emph{narrower/wider} than the parent ($|a|$). They translate fluently among
    an equation, a table, and a pre-drawn graph, reading domain, range, intercepts, and
    increasing/decreasing intervals off the V. Lesson~2.2 \emph{solved} $|ax+b|=c$; today that same
    expression becomes a \emph{graph}, and the two solutions are exactly where the V meets the line
    $y=c$. This gives an early, informal look at \textbf{vertex, axis of symmetry, and min/max} that
    Unit~3 formalizes for parabolas.
    \\[2pt]
    \textbf{Standards (2023 VA SOL):} This lesson applies the Algebra~2 \emph{transformation lens}
    \textbf{A2.F.1} --- write and graph a function using transformations of a parent function
    $f(x)+k$, $f(x+k)$, $f(kx)$, and $kf(x)$ (\textbf{A2.F.1b}, \textbf{A2.F.1c}) --- to the absolute
    value parent as the entry example, and reads the resulting \emph{characteristics} with
    \textbf{A2.F.2}: domain, range, zeros, and intercepts (\textbf{A2.F.2a}), increasing/decreasing
    intervals (\textbf{A2.F.2c}), absolute (global) minimum/maximum (\textbf{A2.F.2d}), and evaluating
    $f(x)$ from a graph or equation (\textbf{A2.F.2f}). Absolute value is itself a
    \emph{piecewise-defined} function --- one of the families A2.F.2 names --- which motivates
    Lesson~2.4.
\end{tcolorbox}

\renewcommand{\arraystretch}{1.4}
\begin{skillbox}[Priority Ideas \& Skills]{goldbox}
\begin{tabularx}{\linewidth}{@{}X|X@{}}
\textbf{Priority Ideas \& Skills} & \textbf{Key Understandings (the ``why'')} \\[2pt]
\hline
\vspace{-6pt}
\begin{itemize}[leftmargin=1.1em, nosep, after=\vspace{2pt}]
  \item Graph the parent $f(x)=|x|$ from a table and read its features (V, vertex, axis, min).
  \item From $g(x)=a|x-h|+k$, name the \textbf{vertex} $(h,k)$ and \textbf{axis} $x=h$.
  \item Decide \textbf{opens up/down} from the sign of $a$, and \textbf{min vs.\ max} value $k$.
  \item Read the effect of $|a|$: \textbf{narrower} ($|a|>1$) or \textbf{wider} ($|a|<1$).
  \item Describe a graph as \textbf{shifts, a stretch, and a reflection} of $y=|x|$.
  \item Read domain, range, intercepts, and increasing/decreasing intervals off the V.
\end{itemize}
&
\vspace{-6pt}
\begin{itemize}[leftmargin=1.1em, nosep, after=\vspace{2pt}]
  \item $|x|$ is a \emph{distance}, so its graph can't dip below its vertex --- that makes the V and its min/max.
  \item Inside the bars acts \emph{horizontally} and \emph{backwards}: $|x-h|$ moves the vertex to $x=h$.
  \item Outside the bars acts \emph{vertically} the way you'd expect: $+k$ up, $a$ scales, $a<0$ flips.
  \item $a$ is the slope of the right-hand ray --- it controls steepness and which way the V opens.
  \item Every absolute value graph is the \emph{one} parent V after moving, scaling, and possibly flipping.
  \item The vertex is the graph's \emph{turning point}: where increasing meets decreasing (previewing Unit~3).
\end{itemize}
\end{tabularx}
\end{skillbox}

\begin{skillbox}[{Vocabulary, Concepts, \& Theorems}]{sky}
\renewcommand{\arraystretch}{1.3}
\begin{tabularx}{\linewidth}{@{}l X@{}}
\textbf{Absolute value parent} & $f(x)=|x|$: a V opening up, vertex at the origin, made of the rays $y=x$ ($x\ge 0$) and $y=-x$ ($x<0$). \\
\textbf{Vertex} & The turning point of the V --- its lowest point (opens up) or highest point (opens down). \\
\textbf{Axis of symmetry} & The vertical line $x=h$ through the vertex; the two rays are mirror images across it. \\
\textbf{Vertex form} & $g(x)=a|x-h|+k$, with vertex $(h,k)$ and axis $x=h$. \\
\textbf{Reflection} & $a<0$ flips the V to open \emph{downward} (over a horizontal line). \\
\textbf{Vertical stretch / compression} & $|a|>1$ makes the V \emph{narrower}; $0<|a|<1$ makes it \emph{wider}. \\
\textbf{Minimum / maximum} & The $y$-value $k$ at the vertex: a \emph{minimum} if $a>0$, a \emph{maximum} if $a<0$. \\
\end{tabularx}
\end{skillbox}

\begin{fixedskillbox}[Activate Prior Knowledge \& Spiral Review]{forestbg}
    The Warm-Up rehearses the three prerequisites today leans on: \textbf{(1)} evaluate absolute
    values and build a table for $y=|x|$, so students \emph{see} why the graph turns at the origin
    and never dips below it; \textbf{(2)} recall the transformation language from Lesson~2.1 ---
    ``$+k$'' shifts up/down and a negative coefficient \emph{reflects} --- the exact moves we now
    apply to a new parent; and \textbf{(3)} read the vertex and axis of symmetry off a pre-drawn V.
    Debrief item~1 aloud: because $|x|$ is a \emph{distance}, the two sides of the table mirror each
    other --- that symmetry \emph{is} the axis of symmetry.
\end{fixedskillbox}

\begin{skillbox}[Hook]{forestbg}
    You drive a straight highway toward a rest stop, reach it after $2$ hours, and keep going past it.
    Your \textbf{distance from the rest stop} (in tens of miles) is $d(t)=30\,|t-2|$. Ask:
    \textbf{``How far are you at the start? At $t=2$? At $t=4$?''} ($60$, $0$, $60$.) \textbf{``When
    are you closest, and what is that distance?''} (at $t=2$, distance $0$). Plot the three points and
    the shape is a \textbf{V} with its lowest point at $(2,0)$. Land the idea: a distance is never
    negative, so its graph \emph{bottoms out} at a single turning point --- the \textbf{vertex} --- and
    is symmetric about the vertical line through it. Every absolute value function is this same V,
    moved and stretched. Today we learn to read the move straight off the equation.
\end{skillbox}

\begin{skillbox}[Lesson]{forestbg}
\begin{multicols}{2}
\textbf{1. The parent $f(x)=|x|$.} Build the table $x=-2,-1,0,1,2 \to y=2,1,0,1,2$. The graph is a
\textbf{V} opening up with \textbf{vertex} $(0,0)$, \textbf{axis} $x=0$, \textbf{minimum} $0$. Domain
all reals; range $y\ge 0$. It \emph{decreases} on $(-\infty,0]$ and \emph{increases} on $[0,\infty)$.
Note it is the two lines $y=x$ and $y=-x$ pasted at the vertex --- a preview of piecewise (2.4).

\textbf{2. Up/down and left/right.} \textbf{Outside} the bars behaves normally: $y=|x|+k$ shifts the
whole V up ($k>0$) or down. \textbf{Inside} the bars behaves \emph{horizontally and backwards}:
$y=|x-h|$ moves the vertex to $x=h$ (\emph{right} for $-h$, left for $+h$). Anchor with $y=|x-3|$
(vertex $(3,0)$) and $y=|x+2|$ (vertex $(-2,0)$).

\textbf{3. Stretch, compress, reflect.} $y=a|x|$ scales the V: $|a|>1$ \textbf{narrower},
$0<|a|<1$ \textbf{wider}. A \emph{negative} $a$ \textbf{reflects} it to open \emph{downward}. Compare
$y=2|x|$, $y=\tfrac12|x|$, $y=-|x|$. Stress that $a$ is the slope of the right-hand ray.

\columnbreak

\textbf{4. Vertex form $g(x)=a|x-h|+k$.} All three moves at once. \textbf{Vertex} $(h,k)$;
\textbf{axis} $x=h$; \textbf{opens up} if $a>0$ (so $k$ is a \emph{minimum}), \textbf{down} if $a<0$
(so $k$ is a \emph{maximum}); $|a|$ sets the width. Model $g(x)=-2|x-1|+3$: vertex $(1,3)$, axis
$x=1$, opens \emph{down}, \emph{maximum} $3$, narrow. Read $g(0)=-2|{-1}|+3=1$ straight from the rule.

\textbf{5. Equation $\leftrightarrow$ graph $\leftrightarrow$ table.} Practice all three directions,
but \emph{never sketch from scratch}: match an equation to the correct pre-drawn V, read a vertex off
a graph, or complete a table and name the vertex. Connect back to 2.2: the solutions of $|x-1|=2$ are
the two $x$-values where $y=|x-1|$ crosses the line $y=2$.
\end{multicols}
\end{skillbox}

\begin{skillbox}[Active Monitoring]{redbox}
Circulate during the notes practice and Tier work. Watch for and correct:
\begin{itemize}[leftmargin=1.2em, nosep]
  \item \textbf{horizontal-shift sign flip}: reading $|x-3|$ as a shift \emph{left} $3$ instead of right;
  \item confusing inside vs.\ outside the bars: applying $+k$ horizontally or the shift vertically;
  \item forgetting that $a<0$ opens the V \emph{down}, so $k$ becomes a \textbf{maximum}, not a minimum;
  \item calling a wider graph ``$|a|>1$'': $|a|>1$ is \emph{narrower}, $0<|a|<1$ is wider;
  \item naming the vertex as $(-h,k)$ or $(h,-k)$ --- mis-reading the signs in $a|x-h|+k$;
  \item saying the range is ``all reals'' --- the V stops at $k$ (range $y\ge k$ up, $y\le k$ down);
  \item swapping increasing/decreasing intervals, or splitting them at $0$ instead of at $x=h$.
\end{itemize}
Cold-call prompts: ``Where is the vertex, and how do you see it in the equation?'' ``Does it open up
or down --- how do you know?'' ``Is $k$ a max or a min here?'' ``Narrower or wider than the parent?''
\end{skillbox}

\begin{skillbox}[Group Work \& Differentiation]{redbox}
\textbf{Reading \& Building V's} activity (tiered; mirrors the notes).
\begin{multicols}{3}
\textbf{Tier R --- Remediate.} Complete tables for $y=|x|$ and one shift; read the vertex, axis, and
minimum off pre-drawn V-graphs; match two simple equations to their graphs.

\columnbreak
\textbf{Tier A --- Approaching.} From vertex-form equations, state vertex, axis, opens up/down, and
min/max; match equations to graphs; describe each as transformations of $y=|x|$; evaluate $g(x)$ at a
point.

\columnbreak
\textbf{Tier E --- Extension.} Model a distance/time context with an absolute value function, write
its equation, and interpret the vertex; write the equation of a V from a described graph; reason about
how the sign of $a$ forces a max vs.\ a min.
\end{multicols}
\end{skillbox}

\begin{skillbox}[Individual Work \& Assessment]{redbox}
\textbf{Exit Ticket} (independent, no notes): read the vertex, axis, range, and an increasing interval
off a pre-drawn V; from $g(x)=-2|x-3|+5$ name the vertex, axis, direction of opening, and the
max/min value; and one \textbf{SOL-style multiple-choice} item (describe the transformation from
$y=|x|$ to $y=|x+4|-2$). Collect and sort into ``got it / shift-sign slip / up-vs-down slip /
vertex-read slip / range slip'' to decide whether Lesson~2.4 opens with a quick re-teach of vertex
form.
\end{skillbox}

\begin{skillbox}[Reinforcement \& Extension]{goldbox}
\textbf{Homework:} recall the parent's characteristics; match equations to pre-drawn V-graphs; for a
set of vertex-form equations state vertex, axis, opens up/down, and max/min; describe transformations
from the parent; complete a table and name a vertex; and model a distance context. \textbf{Extension:}
write the equation of a V from its graph, and reason about when the vertex is a maximum.
\textbf{Preview:} Lesson~2.4, \emph{Piecewise-Defined Functions} --- today's V is really \emph{two}
lines pasted at the vertex ($y=-x$ then $y=x$); next we write functions defined by different rules on
different pieces of the domain, with the greatest-integer/step function as a second example.
\end{skillbox}
% ── Teacher notes (migrated from the component keys) ──
% One per component, titled for it. They live here rather than in the _key
% files so that every component is the same length blank and keyed.

\begin{teachernote}[Warm-Up]
Item~1 is the seed of the whole lesson: the table is \emph{symmetric} because $|x|$ measures distance,
so $x$ and $-x$ give the same $y$ --- that symmetry becomes the axis of symmetry, and the single
smallest value ($0$ at $x=0$) becomes the vertex/minimum. Item~2 reactivates the Lesson~2.1 language
($+k$ shifts, a negative reflects) that we now aim at the new parent $y=|x|$. In item~3, tie the
dashed line $x=0$ to the fold line of the V --- the two rays are mirror images across it.
\end{teachernote}

\begin{teachernote}[Guided Notes]
The spine is the distance meaning: a distance can't go below the vertex, which creates the single
turning point (min/max) and the mirror symmetry (axis). Two errors to pre-empt: (1) the horizontal
\emph{sign flip} --- $|x-3|$ moves \emph{right}, because $x=3$ is where the inside equals zero; have
students set ``inside $=0$'' to locate the vertex every time. (2) The $a<0$ case --- insist that
opening \emph{down} turns $k$ into a \emph{maximum}, and that the range flips to $y\le k$. For width,
tie ``narrower'' to $|a|>1$ (steeper rays) and ``wider'' to $0<|a|<1$. In Section~4, model reading
$g(0)$ \emph{straight from the equation} so students see the function as a machine, not just a
picture. Connect to Lesson~2.2: the solutions of $|x-1|=2$ are the two $x$'s where this V meets the
line $y=2$.
\end{teachernote}

\begin{teachernote}[Group Activity]
Tier~A item~2 and Tier~E items~2--3 are the conceptual core --- push students to locate the vertex by
setting the inside of the bars to zero, and to read $a$ as ``sign $=$ direction, size $=$ width.''
Common slips: in A1(b), forgetting that $a<0$ makes $0$ a \emph{maximum}; in A3, reading $|x-1|$ as a
left shift; in E2, computing $a$ with the wrong run or dropping the negative (the graph opens down, so
$a<0$). For E1, stress that the vertex $(4,0)$ is the \emph{closest} moment because a distance can't be
negative --- the same distance idea that makes every V.
\end{teachernote}

\begin{teachernote}[Exit Ticket]
Item~1 checks reading the vertex, axis, range, and an increasing interval off a V whose vertex is
\emph{below} the $x$-axis --- a good test of ``range $=y\ge k$'' and ``increasing to the right of the
vertex.'' Item~2 checks vertex form with $a<0$: students who say ``min'' missed that opening down makes
$k=5$ a \emph{maximum}. Item~3 separates students who apply the horizontal sign flip correctly
(inside, backwards) from those who read $+4$ as a right shift. Sort into ``got it / shift-sign slip /
up-vs-down slip / vertex-read slip / range slip'' to decide whether Lesson~2.4 opens with a quick
re-teach of vertex form.
\end{teachernote}

\begin{teachernote}[Homework]
Watch item~4(b): $-3|x+1|$ bundles \emph{three} moves (left $1$, reflect, stretch $3$) --- students
often report only the shift. Item~3(d) tests $a<0$ with a fraction: the graph is \emph{wider} and opens
down, so $2$ is a \emph{maximum}. In item~6, tie the vertex $(6,0)$ to ``closest'' because a distance
can't be negative. Extension~(a): $a$ is the slope of the right ray --- rise $4$ over run $2$ gives
$a=2$; a common slip is to forget that a point $2$ units right of the vertex rose $4$, not $1$.
\end{teachernote}

\end{document}
